\subsection{D. The letters and the printed EGA statements}

We distinguish a false assertion, an invalid proof of a true assertion,
an unnecessary hypothesis, and a convention.  In particular the letter
of 29 October concerns \emph{EGA IV} 19.5.5 and 19.5.6, not
\emph{EGA $0_{\rm IV}$}; the latter numbers denote different text.
The primary text of $0_{\rm IV}$ used below is the 1964 printing,
pp.~85--86, 94, and 124--127, opened both as a scan and as searchable
text.  The letters are cote 23, scans 044--065.
\cite{algEGAIV1,algEGAIV4}

\subsubsection{D.1. A repair of \texorpdfstring{$0_{\rm IV}$}{0-IV} 19.3.12}

\begin{lemma}[Proved: the linear compactness needed here]\label{SA23-D-1}
Let $R$ be a ring and let $(M_i)_{i\in I}$ be Artinian $R$-modules.
Consider any set of affine $R$-linear equations in variables
$z_i\in M_i$, each equation involving only finitely many variables
and taking values in an $R$-module.  If every finite subsystem has a
solution, then the entire system has a solution.
\end{lemma}
\begin{proof}
A family of cosets of submodules of one Artinian module, with the
finite intersection property, has nonempty intersection.  Indeed,
among the submodules underlying its nonempty finite intersections
choose a minimal one, using the descending chain condition.  Its
coset is unchanged on intersection with any further member of the
family, since two nonempty cosets of the same submodule, one contained
in the other, are equal.

Well-order the variables.  We choose their values successively while
preserving finite solvability of the original equations together with
the equations fixing all values already chosen.  For the next
variable $z_i$, every finite subsystem projects to a nonempty coset
of a submodule of $M_i$: the solution set of a finite affine linear
system is a coset, and so is its image under a linear projection.
These cosets have the finite intersection property by the induction
hypothesis.  Choose a value in their intersection by the preceding
paragraph.  At a limit ordinal finite solvability persists, since a
finite subsystem mentions only finitely many previous choices.  At
the end all the equations hold.  This proof does not assert any
analogous compactness for polynomial equations.
\end{proof}

\begin{lemma}[Proved: closed quotients of a pro-Artinian ring]\label{SA23-D-2}
Let $C$ be a topological ring with a directed fundamental system of
open ideals $(U_\lambda)$ such that $C_\lambda=C/U_\lambda$ is
Artinian and $C\simeq\varprojlim_\lambda C_\lambda$ topologically.
For every closed ideal $D$ of $C$, the canonical map
\[
 C/D\longrightarrow\varprojlim_\lambda C/(D+U_\lambda)
\]
is an isomorphism of topological rings.  In particular the closed
quotient is again complete and pro-Artinian.
\end{lemma}
\begin{proof}
Injectivity is the identity
$D=\bigcap_\lambda(D+U_\lambda)$, which is the definition of
closedness for this topology.  Given compatible residues
$d_\lambda\in C/(D+U_\lambda)$, use variables
$c_\lambda\in C_\lambda$ with the linear equations prescribing
their images to be $d_\lambda$ and requiring compatibility under the
quotient maps.  Each $C_\lambda$ is an Artinian $C$-module.  A finite
subsystem is solved by choosing an index dominating all its indices
and lifting the prescribed residue there to an element of $C$.
Lemma~\ref{SA23-D-1} gives compatible $c_\lambda$, hence an element
of $C$.  This proves surjectivity.  Both topologies have as a
fundamental system of open ideals the kernels of the displayed maps
to $C/(D+U_\lambda)$, proving the topological assertion.
\end{proof}

\begin{theorem}[Proved: the statement of $0_{\rm IV}$ 19.3.12 is true]\label{SA23-D-3}
Let $A$ be a topological ring, $B$ a formally smooth topological
$A$-algebra, and $C$ a topological $A$-algebra.  Suppose that
$C=\varprojlim C/U_\lambda$ topologically with each $C/U_\lambda$
Artinian.  Let $J\subset C$ be a closed ideal whose image in every
$C/U_\lambda$ is nilpotent; this follows in particular when $J$ is
topologically nilpotent as in the printed proposition.  Suppose also
that the square of every open ideal of $B$ is open.  Every continuous
$A$-algebra homomorphism $u:B\to C/J$ lifts to a continuous
$A$-algebra homomorphism $v:B\to C$.
\end{theorem}
\begin{proof}
We first prove the assertion when $J^2=0$.  Write
$J_\lambda=(J+U_\lambda)/U_\lambda$ and
$u_\lambda:B\to C_\lambda/J_\lambda$ for the induced map.
Formal smoothness gives an $A$-algebra lift
$w^0_\lambda:B\to C_\lambda$ for each $\lambda$.
The possible algebra lifts have the form
$w_\lambda=w^0_\lambda+z_\lambda$, where the variables
$z_\lambda(b)\in J_\lambda$ for $b\in B$ satisfy the equations
\[
\begin{split}
 z_\lambda(b+b')&=z_\lambda(b)+z_\lambda(b'),\\
 z_\lambda(bb')&=w^0_\lambda(b)z_\lambda(b')
                  +w^0_\lambda(b')z_\lambda(b),\\
 z_\lambda(a)&=0\quad(a\in A),\qquad z_\lambda(1)=0.
\end{split}
\]
Here $a$ means its image in $B$.  These are linear equations over
$C$ in Artinian $C$-modules $J_\lambda$.  Their validity is exactly
where the square-zero hypothesis is used.  The compatibility
equations $q_{\lambda\mu}w_\mu=w_\lambda$ are affine linear of
the same kind.  Every finite subsystem has a solution: choose an
index $\nu$ dominating the finitely many indices and lift
$u_\nu$ by formal smoothness, then project this lift to those
indices.  Lemma~\ref{SA23-D-1} supplies compatible lifts, hence an
$A$-algebra map $v:B\to C$ reducing to $u$; reduction is justified
by Lemma~\ref{SA23-D-2} with $D=J$.

This algebra map is continuous.  For each $\lambda$, continuity of
$u_\lambda$ gives an open ideal $V\subset B$ with
$u_\lambda(V)=0$.  Its image under $w_\lambda$ lies in
$J_\lambda$, so $w_\lambda(V^2)=0$.  The ideal $V^2$ is open by
hypothesis.  Thus every $w_\lambda$ is continuous, and so is $v$.

For general $J$, put $J_n=\overline{J^n}$ for $n\geq1$; in
particular $J_1=J$.  Continuous multiplication gives
$J_n^2\subset\overline{J^{2n}}\subset J_{n+1}$.  By
Lemma~\ref{SA23-D-2}, each $C/J_n$ is complete and pro-Artinian.
The map $C/J_{n+1}\to C/J_n$ has closed square-zero kernel
$J_n/J_{n+1}$.  Applying the square-zero case repeatedly constructs
compatible continuous lifts $u_n:B\to C/J_n$ with $u_1=u$.

We finally verify the comparison used to take their limit.  For
each $\lambda$, nilpotence of the image of $J$ implies
$J^{n(\lambda)}\subset U_\lambda$, hence
$J_{n(\lambda)}\subset U_\lambda$ since $U_\lambda$ is closed.
Thus $\bigcap_nJ_n=0$ and
\[
 C\longrightarrow\varprojlim_{n\geq1}C/J_n
\]
is injective.  To show surjectivity choose representatives $c_n$
of compatible classes.  For $m\geq n$ their difference lies in
$J_n$.  Modulo each $U_\lambda$ the sequence is therefore
eventually constant; these eventual residues form an element of
$\varprojlim C_\lambda=C$.  Its difference from $c_n$ belongs
to $J_n$ by closedness.  This proves surjectivity.  The comparison
is a homeomorphism: any open $U_\lambda$ contains some $J_n$,
and $U_\lambda/J_n$ is open in $C/J_n$.  The inverse limit of
the continuous maps $u_n$ consequently gives the required
continuous map $v:B\to C$.
\end{proof}

\begin{proposition}[Refuted proof step; no counterexample to the theorem]\label{SA23-D-4}
The proof printed in $0_{\rm IV}$ 19.3.12, pp.~85--86, wrongly
treats sets of lifts through arbitrary nilpotent ideals as affine
linear varieties.  Deligne's polynomial compactness counterexample
on cote 23, scan 060, is correct, but it does not refute
Theorem~\ref{SA23-D-3}.
The final Err$_{\rm IV}$~8 explicitly deletes the proposition
because its proof is erroneous; it identifies the purported
linear varieties as the problem.  It does not supply a
counterexample to the statement.
\cite[Err$_{\rm IV}$~8, p.~347]{algEGAIV4}
\end{proposition}
\begin{proof}
For an infinite field $k$ consider variables $x$ and $y_a$ for
$a\in k$, with equations $(x-a)y_a=1$.  Any finite collection is
soluble: choose $x$ different from the finitely many specified
$a$ and set $y_a=(x-a)^{-1}$.  A simultaneous solution would
require $x\ne a$ for every $a\in k$, which is impossible.
This is a counterexample to the polynomial compactness principle
used by the printed proof.  In a lift $w^0+z$ the multiplicativity
equation contains $z(b)z(b')$ when the kernel is not square-zero,
so it need not be affine linear.  The repaired proof applies linear
compactness only to square-zero kernels, then takes the countable
tower of closed powers.  The statement therefore remains valid.
The literature searches recorded in the source log found the printed
proposition and its mirrors, but no prior source for this particular
repair; we make no priority claim.
\end{proof}

\subsubsection{D.2. The filtered and graded assertions}

\begin{proposition}[Proved: IV 19.5.5, the implication in the letter]\label{SA23-D-5}
Let $R$ be a ring, $M=F_0M\supset F_1M\supset\cdots$ a filtered
$R$-module, $f=(f_1,\ldots,f_r)$, and $I=(f)$.  Put
$G_nM=\sum_{i=0}^n I^iF_{n-i}M$ and $X=\operatorname{gr}_F M$.
If the canonical graded surjection
\[
 \psi:(X/IX)[T_1,\ldots,T_r]\longrightarrow
       \operatorname{gr}_G M,
 \qquad [m]T^\alpha\longmapsto[f^\alpha m]
\]
is an isomorphism, then $f$ is $X$-quasi-regular.  Here $[m]\in X_p$
and $T^\alpha$ has degree $|\alpha|$, so the target degree is
$p+|\alpha|$.  No separatedness assumption is needed.
This is the final printed IV 19.5.5(b)$\Rightarrow$(c), pp.~206--209;
the footnote on p.~207 explicitly credits Deligne's proof.  The
corresponding numbers in $0_{\rm IV}$, p.~94, are a different remark
and a proof of Theorem 19.5.3, not the claims criticized in this letter.
\end{proposition}
\begin{proof}
Injectivity of $\psi$ has this useful consequence: if a finite sum
\[
 \sum_{p+|\alpha|=n} f^\alpha m_{p,\alpha}\in G_{n+1}M,
 \qquad m_{p,\alpha}\in F_pM,
\]
then every $m_{p,\alpha}$ lies in $F_{p+1}M+IF_pM$.
This is simply the direct-sum decomposition of the degree-$n$
part of the source of $\psi$.

We prove by induction on $q\geq1$ that
\[
 I^qF_pM\cap F_{p+1}M
 \subset\sum_{i=1}^q I^{q-i}F_{p+i}M
 \quad\hbox{for every }p\geq0.                         \tag{D.5.1}
\]
The case $q=1$ is immediate.  For
$m\in I^{q+1}F_pM\cap F_{p+1}M$, the induction hypothesis writes
$m=\sum_{i=1}^q\sum_{|\alpha|=q-i}f^\alpha m_{i,\alpha}$,
where $m_{i,\alpha}\in F_{p+i}M$.  But
$m\in I^{q+1}F_pM\subset G_{p+q+1}M$.  The displayed consequence
of injectivity, in total degree $p+q$, gives
$m_{i,\alpha}\in F_{p+i+1}M+IF_{p+i}M$ for every $i,\alpha$.
Substitution yields (D.5.1) for $q+1$.

To prove quasi-regularity take $x_\alpha\in F_pM$, $|\alpha|=q$,
such that
$\sum f^\alpha x_\alpha\in F_{p+1}M+I^{q+1}F_pM$.
Subtracting suitable elements of $IF_pM$ from the $x_\alpha$
makes this sum belong to $F_{p+1}M$ without changing their classes
in $F_pM/(F_{p+1}M+IF_pM)$.  Formula (D.5.1) then expresses the
sum using terms $f^\beta y_{i,\beta}$ with
$y_{i,\beta}\in F_{p+i}M$ and $|\beta|=q-i$, $i>0$.
The resulting equality, in total degree $p+q$, and injectivity of
$\psi$ give $x_\alpha\in F_{p+1}M+IF_pM$.  This is exactly
injectivity of
$(X/IX)[T]_q\to I^qX/I^{q+1}X$ on each $F$-degree.  Surjectivity
is automatic from the generators of $I$, proving the assertion.
\end{proof}

\begin{example}[Refuted: the final IV 19.5.6(ii); corrected scope]\label{SA23-D-6}
The final 1967 text, p.~209, assumes only that the $f_i$ have no
homogeneous component of degree zero.  Under that wording its
quasi-regular$\Rightarrow$regular assertion is false.  For any field
$k$, take the ordinarily graded ring and module
\[
 R=M=k[x,y,z],\qquad f=((1-x)y,(1-x)z,x).
\]
Every $f_i$ has strictly positive degree components, $M$ is bounded
below, and $f$ is Koszul-regular and quasi-regular.  It is not an
ordinary regular sequence.  A valid corrected statement is obtained
by requiring the $f_i$ to be homogeneous of positive degree; more
generally the exact sufficient hypothesis is separatedness of each
successive quotient for the $I$-adic topology, as in
$0_{\rm IV}$ 15.1.9.
\end{example}
\begin{proof}
The sequence generates $(x,y,z)$, since
$y=(1-x)y+xy$ and $z=(1-x)z+xz$.  The latter ideal is generated
by the ordinary regular sequence $(x,y,z)$.  The generator-invariance
and quasi-regularity results of
Theorems~\ref{SA23-C-3} and \ref{SA23-C-2} show that $f$ is
Koszul-regular and quasi-regular.  Its first element is a
nonzerodivisor, because $R$ is a domain.  In
$R/((1-x)y)$, however, the nonzero class of $y$ is killed by
$(1-x)z$, since $y(1-x)z=z(1-x)y$.
The class of $y$ is nonzero: cancellation in the domain would
otherwise make $1-x$ a unit in $k[x,y,z]$.  Thus the second
successive injectivity condition fails.

For completeness, the separation criterion needed for the correction
follows directly from quasi-regularity.  If
$(M/IM)[T]\simeq\operatorname{gr}_I M$, then multiplication by
$f_1$ is strict and injective whenever $M$ is $I$-adically separated:
if $f_1m\in I^{n+1}M$, comparison of its initial form shows
$m\in I^nM$ by induction on its filtration order; for
$f_1m=0$ this puts $m$ in every $I^nM$.  Strictness means
$f_1M\cap I^nM=f_1 I^{n-1}M$ for $n\geq1$.
It identifies the associated graded of $M/f_1M$ with
$(M/IM)[T_2,\ldots,T_r]$.  Repeating this argument, assuming
separatedness of the successive quotients, proves regularity.
When the $f_i$ are homogeneous of positive degree, these quotients
remain graded and bounded below; $I^n$ raises the least possible
degree by at least $n$, proving the required separatedness.
No finite generation of the graded module is used.  If one's
definition of regularity requires a nonzero final quotient, it is
included as a separate hypothesis (and holds in the example).

The earlier, stronger claim that \emph{every} subquotient is separated
also fails even with homogeneous generators: for $R=k[T]$ and
$I=(T)$, the quotient $R/(T-1)$ is nonzero and equals its own
$I$-multiple.  Likewise the relative Noetherian finiteness mentioned
in Deligne's December letter is insufficient for arbitrary
$A$-subquotients: take $A=k[x]_{(x)}$, $B=A[[T]]$, $M=B$.
The $A$-submodule $A[T]\subset B$ has quotient
$A[T]/(xT-1)\simeq k(x)$, on which $x$ is invertible.  Restricting
to finite $B$-subquotients and $IB\subset\operatorname{Jac}(B)$
does give separatedness by the Noetherian Krull intersection theorem.
\cite[Lemma 10.51.5(2)]{algStacksKrull}
\end{proof}

\subsubsection{D.3. Differentials, infinitesimal neighborhoods, and completion}

\begin{proposition}[Proved: the printed sign correction to $0_{\rm IV}$ 20.4.6]\label{SA23-D-7}
Let $A\to B$ be a ring map, $D=B\otimes_A B$, and
$I=\ker(D\xrightarrow{\mu}B)$.  With $B$ acting through
$p_1(b)=b\otimes1$, the convention compatible with the usual
first-order increment is
\[
 d(b)=p_2(b)-p_1(b)=1\otimes b-b\otimes1\pmod{I^2}.
\]
It makes $(I/I^2,d)$ the universal differential module.  The
opposite sign originally printed in $0_{\rm IV}$ 20.4.6 is also
abstractly universal, but the final printed errata do change this
specified convention.  Their exact wording has a further defect:
the first clause of Err$_{\rm IV}$~12 says to replace $p_2-p_1$
by $p_1-p_2$, although the original 20.4.6.1 already has
$p_1-p_2$.  That clause has its old and new expressions reversed.
Its next clause unambiguously replaces $x\otimes1-1\otimes x$
by $1\otimes x-x\otimes1$; Err$_{\rm IV}$~13 likewise changes
the relation in 20.4.8 to $p_1=p_2-d$.
These latter two corrections both specify $d=p_2-p_1$.
Thus Deligne's requested sign correction is represented in the
final errata, but their first substitution itself requires the
correction just identified.\cite[Err$_{\rm IV}$~12--13,
pp.~347--348]{algEGAIV4}  None of this makes the original
abstract universal property false.
\end{proposition}
\begin{proof}
For an $A$-derivation $\delta:B\to M$, where $M$ is a $B$-module,
define $\theta:D\to M$ by
$\theta(b\otimes c)=b\delta(c)$.  The Leibniz identity gives
$\theta(uv)=\mu(u)\theta(v)+\mu(v)\theta(u)$, so $\theta(I^2)=0$.
Its restriction induces a $B$-linear map $I/I^2\to M$ taking
$d(b)$ to $\delta(b)$.  The elements $d(b)$ generate $I/I^2$:
if $\sum b_i c_i=0$, then
$\sum b_i\otimes c_i=\sum (b_i\otimes1)(1\otimes c_i-c_i\otimes1)$.
Thus the map is unique.  Replacing $d$ by $-d$ and the representing
map by its negative gives the opposite convention.  This proves
universality in either sign and identifies the required conversion.

The canonical identification in the surrounding printed text fixes
which sign appears.  Write $P=D/I^2$, $\Omega=I/I^2$, and
$\epsilon:P\to B$ for the augmentation.  The isomorphism in
$0_{\rm IV}$ 20.4.8(i), which is required to be the identity on
$\Omega$, is
\[
 \Phi:P\longrightarrow B\ltimes\Omega,\qquad
 z\longmapsto(\epsilon z,z-p_1(\epsilon z)).
\]
It satisfies $\Phi p_1(b)=(b,0)$ and
$\Phi p_2(b)=(b,p_2(b)-p_1(b))=(b,d(b))$.  With the original
$d_{\rm old}=p_1-p_2$, the latter formula is instead
$(b,-d_{\rm old}(b))$.  This explicit computation checks the fixed
canonical identification, rather than discarding its sign.  The old
proof writes $p_2=p_1-d_{\rm old}$ and thereby still establishes a
bijection of homomorphisms with derivations; the final correction
aligns the named differential with the second component of $\Phi p_2$.
\cite[20.4.1--20.4.8, pp.~123--126]{algEGAIV1}
\end{proof}

\begin{proposition}[Proved: the normality hypothesis in 20.4.13(iv) is redundant]\label{SA23-D-8}
For an inclusion of domains $A\subset B$, put
$K=\operatorname{Frac}(A)$ and $L=\operatorname{Frac}(B)$.
There is a canonical isomorphism
\[
 L\otimes_B\Omega_{B/A}\simeq\Omega_{L/K}.
\]
Consequently $\Omega_{B/A}$ is a torsion $B$-module if and only if
$\Omega_{L/K}=0$.  In particular it is torsion whenever $L/K$ is
algebraic and separable.  The printed $0_{\rm IV}$ 20.4.13(iv)
is true, but neither integral closedness of $A$ nor integrality of
$B/A$ is needed beyond the field-extension condition just stated.
Err$_{\rm IV}$~14 already incorporates the removal of normality
and supplies the cleared-minimal-polynomial argument.
\cite[Err$_{\rm IV}$~14, p.~348]{algEGAIV4}
\end{proposition}
\begin{proof}
Every $A$-derivation from $B$ to an $L$-vector space extends uniquely
to $L$ by $d(b/s)=(s\,db-b\,ds)/s^2$.  It kills all fractions
from $A$, hence is a $K$-derivation.  This proves the canonical
isomorphism by the representing properties.  Localization at all
nonzero elements of a domain kills a module exactly when each
element is killed by some nonzero scalar, which proves the stated
equivalence.  If $a\in L$ is separable algebraic over $K$, its
minimal polynomial $P$ has $P'(a)\ne0$; differentiating $P(a)=0$
gives $da=0$.  Applying this to every $a$ gives
$\Omega_{L/K}=0$.  We do not assert the converse that zero
differentials forces an algebraic separable extension in arbitrary
positive characteristic.
\end{proof}

\begin{proposition}[Proved: the final IV 16.1.5 and 16.1.6]\label{SA23-D-9}
For an immersion $f:Y\to X$ of schemes, the associated graded
normal algebra is quasi-coherent on $Y$, and every infinitesimal
neighborhood $Y^{(n)}$ is a scheme immersed in $X$.  A sheaf-of-rings
section $\mathcal O_Y\to\mathcal O_{Y^{(n)}}$ of the augmentation
makes the latter and its lower-order quotients quasi-coherent
$\mathcal O_Y$-algebras; it induces the canonical module structures
on the associated graded pieces.  If $f$ is locally of finite
presentation, every normal graded piece is of finite type.

These are the assertions in the final 1967 printing,
\cite[IV 16.1.5--16.1.6, pp.~6--7]{algEGAIV4}, where
\emph{immersion} is explicit.  The disjoint-union map in Deligne's
October letter does not refute this final formulation.
\end{proposition}
\begin{proof}
The assertions are local on an open neighborhood of $Y$ in $X$,
so factor the immersion there as a closed immersion, with
quasi-coherent ideal sheaf $\mathcal I$.  Then $Y^{(n)}$ is the
closed subscheme defined by $\mathcal I^{n+1}$, and the normal
graded pieces are $\mathcal I^q/\mathcal I^{q+1}$ viewed on $Y$.
These descriptions prove the first two assertions and agree on
overlaps, so they glue for the original immersion.

A section of the augmentation gives an $\mathcal O_Y$-module
structure on $\mathcal O_{Y^{(n)}}$.  Its finite filtration by
$\mathcal I^q/\mathcal I^{n+1}$ has graded pieces
$\mathcal I^q/\mathcal I^{q+1}$ with the canonical action:
two lifts of a section of $\mathcal O_Y$ differ by $\mathcal I$,
which annihilates each such quotient.  Starting at $q=n$,
successive extensions of these quasi-coherent modules prove
quasi-coherence of all filtered terms and of
$\mathcal O_{Y^{(n)}}$.  This uses only that quasi-coherent modules
are closed under extensions, a fact checked on each affine open by
the equivalence with modules.

For an immersion locally of finite presentation the local defining
ideal $\mathcal I$ is of finite type.  The surjection
$\operatorname{Sym}^q_{\mathcal O_Y}(\mathcal I/\mathcal I^2)
\twoheadrightarrow\mathcal I^q/\mathcal I^{q+1}$ proves the
finite-type assertion.  In Deligne's example
$Y=\operatorname{Spec}k[T]\amalg\operatorname{Spec}k\to
X=\operatorname{Spec}k[T]$, with the second component sent to the
origin, two distinct points of $Y$ map to that origin.  An immersion
is injective on underlying points, so this example is outside the
final printed hypotheses.  The draft criticism and the final
printed assertion must therefore be distinguished.
\end{proof}

\begin{example}[Refuted: finite presentation does not make the kernel finite]\label{SA23-D-10}
The auxiliary question raised in the letter has a negative answer:
a finitely presented ring map need not have a finitely generated
kernel.  Let $k$ be a field,
$M=\bigoplus_{i\geq1}k$, with $xM=0$, and
$A=k[x]\ltimes M$.  The map
\[
 A\longrightarrow A[1/x]=A[T]/(xT-1)
\]
is finitely presented and has kernel $0\ltimes M$, which is not
finitely generated as an ideal.  For a \emph{surjective} finitely
presented ring map the kernel is finitely generated, as required
in the closed-immersion proof of Proposition~\ref{SA23-D-9}.
\end{example}
\begin{proof}
Every element of $M$ is killed by $x$, while $k[x]$ has no
$x$-power torsion.  Thus localization kills exactly $0\ltimes M$
and identifies the target with $k[x,x^{-1}]$.  The displayed
one-generator, one-relation presentation proves finite
presentation over $A$.  The action on $0\ltimes M$ factors
through $A\to k$, so finitely many ideal generators could span
only a finite-dimensional $k$-subspace of the infinite direct
sum.  Finally, if $A\twoheadrightarrow B$ is finitely presented,
present $B=A[T_1,\ldots,T_s]/(g_1,\ldots,g_t)$ and choose
$a_i\in A$ mapping to the image of $T_i$.  The kernel is generated
by the finitely many substitutions $g_j(a_1,\ldots,a_s)$, by
division by the monic linear polynomials $T_i-a_i$.
\end{proof}

\begin{proposition}[Refuted formula; proved conclusion in III 4.3.6]\label{SA23-D-11}
Let $A$ be a Noetherian local domain whose completion $\widehat A$
is a domain, and set $K=\operatorname{Frac}(A)$ and
$K'=\operatorname{Frac}(\widehat A)$.  The equality
$K'=K\otimes_A\widehat A$, printed in
\cite[III 4.3.6, p.~133, line 7]{algEGAIII1}, is false in general.
The exact replacement is
\[
 K\otimes_A\widehat A=(A\setminus\{0\})^{-1}\widehat A
 \ \subset\ K'.
\]
Equality holds exactly when this localization is a field, or
equivalently when the only prime of $\widehat A$ contracting to
zero in $A$ is zero.  The conclusion that an analytically integral
Noetherian local domain is unibranch remains true.
\end{proposition}
\begin{proof}
For a counterexample take $A=\mathbf Q[x,y]_{(x,y)}$, so
$\widehat A=\mathbf Q[[x,y]]$.  Let $h(x)=\exp(x)-1$ and
$g=y-h(x)$.  Lemma~\ref{SA23-A-03} implies that $h$ is
transcendental over $\mathbf Q(x)$.
Evaluation $y\mapsto h(x)$ gives a surjection
$\widehat A\to\mathbf Q[[x]]$ with kernel $(g)$, by division
by the monic linear polynomial in $y$.  Its contraction to $A$
is zero: a nonzero polynomial vanishing at $h(x)$ would make
$h$ algebraic, and a denominator outside $(x,y)$ evaluates to a
unit.  Therefore $(g)$ is a nonzero prime surviving in
$(A\setminus0)^{-1}\widehat A$, and this localization is not a
field.  Concretely, $1/g\in K'$ cannot equal $b/a$ with
$b\in\widehat A$ and $0\ne a\in A$, since that would put
$a$ in $(g)\cap A=0$.

The corrected formula is the localization tensor identity.  Its
inclusion in $K'$ follows since $\widehat A$ is a domain and
every nonzero element of $A$ remains nonzero by faithful flatness
of completion.  The criterion for equality follows from the
prime correspondence for localization.

To repair the unibranch argument, let $A_\beta\subset K$ be a
finite $A$-subalgebra.  Flatness of $\widehat A$ embeds
$A_\beta\otimes_A\widehat A$ into
$K\otimes_A\widehat A\subset K'$; hence it is a domain.
By finite-module completion it is the maximal-ideal-adic
completion of $A_\beta$.  A finite algebra over a local ring is
semilocal, and its completion is the product of the completions
at its finitely many maximal ideals
\cite[Lemmas 10.97.1 and 10.97.8]{algStacksNoetherianCompletion}.
Every factor is nonzero, so the domain property forces exactly
one maximal ideal.  Thus each $A_\beta$ is local.
The integral closure of $A$ in $K$ is the filtered union of
these finite subalgebras.  Their maximal ideals are compatible
under inclusion, by lying over for integral extensions.  Their
union is consequently a local ring: an element outside the union
of those maximal ideals is already invertible in a finite
subalgebra.  This is precisely unibranchness in III 4.3.6.
\end{proof}

\subsubsection{The meromorphic-section corrections of L\^e D\~ung Tr\'ang}

The following items concern the final printed
\cite[IV 20.1.3--20.1.11, pp.~227--230;
21.1.5--21.1.6, pp.~256--257;
21.2.9, p.~261; 21.6.2, p.~271]{algEGAIV4},
compared with cote 23, scans 045--047.  The notation
$\mathscr S(\mathcal F)$ below always means \emph{locally regular
sections}: a section is required to be a nonzerodivisor at every
stalk of the sheaf of rings $\mathcal F$.
The following table accounts for every uncancelled item on those
three original sheets \cite[scans 045--047]{algManuscript23}.
\begin{center}
\begin{tabular}{ll}
Printed location & Resolution below\\
\hline
20.1.3, section and stalk assertions & \ref{SA23-D-12}, \ref{SA23-D-13}\\
20.1.4 & \ref{SA23-D-14}, localization wording\\
20.1.8 & \ref{SA23-D-14}, the unit assertion\\
20.1.11 & \ref{SA23-D-14}, pullback denominators\\
21.1.5 & \ref{SA23-D-14}, the positive cone\\
21.1.6 & \ref{SA23-D-14}, the divisor inequality\\
21.2.9, after each of (1) and (2) & \ref{SA23-D-15}\\
21.6.2 & \ref{SA23-D-16}
\end{tabular}
\end{center}
On scan 047 the proposed addition to 21.2.12 and the isolated
reference 21.3.6 are visibly crossed out.  They are withdrawn
entries, not two further asserted counterexamples.

\begin{proposition}[Refuted and repaired: IV 20.1.3]\label{SA23-D-12}
For a ringed space $X$, the assignment
\[
 U\longmapsto
 \{\text{nonzerodivisors in }\Gamma(U,\mathcal O_X)\}
\]
need not be a presheaf, even when $X$ is a Noetherian scheme
of finite type over a field.  Its replacement
\[
 \mathscr S(\mathcal O_X)(U)=
 \{s\in\Gamma(U,\mathcal O_X):
 s_x\text{ is a nonzerodivisor in }\mathcal O_{X,x}
 \text{ for every }x\in U\}
\]
is a multiplicative subsheaf.  Equivalently, $s$ belongs to it
if and only if its restriction to every open subset $V\subset U$
is a nonzerodivisor in $\Gamma(V,\mathcal O_X)$.
\end{proposition}
\begin{proof}
For the equivalence, stalkwise injectivity of multiplication by
$s$ implies injectivity on sections over every open.  Conversely,
if $s_xa_x=0$, represent $a_x$ by a section $a$ and shrink its
neighborhood until $sa=0$ there.  The assumed injectivity on
that neighborhood gives $a=0$ and hence $a_x=0$.
The stalkwise condition is local, compatible with restrictions,
and stable under multiplication; gluing of sections proves the
sheaf assertion.

Here is the counterexample to the printed assignment.  Put
$Y=\mathbf P^1_k\times\mathbf A^1_k$, with coordinate $t$
on the second factor, and let $i:\mathbf P^1_k\to Y$ be the
fiber $t=0$.  Set $\mathcal N=i_*\mathcal O_{\mathbf P^1}(-1)$
and
\[
 X=\operatorname{Spec}_Y(\mathcal O_Y\ltimes\mathcal N),
 \qquad
 (a,m)(b,n)=(ab,an+bm).
\]
This is a finite scheme over $Y$, with the same underlying
topological space; it is Noetherian and of finite type over $k$.
We have
$\Gamma(Y,\mathcal O_Y)=k[t]$ and
$\Gamma(Y,\mathcal N)=0$, so
$\Gamma(X,\mathcal O_X)=k[t]$.  Thus $t$ is a nonzerodivisor
globally.  On either standard affine chart of $\mathbf P^1$,
with coordinate $z$, the corresponding open of $X$ has ring
\[
 k[z,t]\ltimes k[z],\qquad t\cdot k[z]=0.
\]
There $t(0,1)=0$ although $(0,1)\ne0$.
Thus restriction does not preserve the printed condition.
For completeness, the two global-section computations follow
on the standard two-chart cover: the common polynomial sections
of $\mathcal O_Y$ are
$k[t,z]\cap k[t,z^{-1}]=k[t]$, whereas a section of
$\mathcal O_{\mathbf P^1}(-1)$ would satisfy
$a(z)=z^{-1}b(z^{-1})$ with both $a,b$ polynomials, forcing
both to be zero.
\end{proof}

\begin{example}[Refuted: the stalk assertion in IV 20.1.3]\label{SA23-D-13}
Even with the corrected sheaf, its stalk
$\mathscr S(\mathcal O_X)_x$ need not consist of all
nonzerodivisors of $\mathcal O_{X,x}$.  Consequently
the meromorphic stalk
$\mathcal M_{X,x}=\mathscr S(\mathcal O_X)_x^{-1}
\mathcal O_{X,x}$ need not be the total quotient ring of
$\mathcal O_{X,x}$.
\end{example}
\begin{proof}
Take an infinite field $k$, put $R=k[x,y]$, and let
\[
 N=\bigoplus_{a\in k}R/(x,y-a),\qquad
 A=R\ltimes N,\qquad X=\operatorname{Spec}A.
\]
The prime $\mathfrak p=(x)\ltimes N$ has local ring
$A_{\mathfrak p}=R_{(x)}$: every summand of $N$ is killed
after localization by $y-a\notin(x)$, and every element of
$N$ has finite support.  Thus $x$ is a nonzero nonunit
nonzerodivisor in $A_{\mathfrak p}$.

We claim
$\mathscr S(\mathcal O_X)_{\mathfrak p}
=A_{\mathfrak p}^{\times}$.
Every unit germ extends to a unit section on a neighborhood,
so one inclusion is immediate.  For the reverse inclusion,
represent any section germ on a principal neighborhood
$D(s)$ of $\mathfrak p$ by $b/s^n$, where
$b,s\in A$ and $s\notin\mathfrak p$.
If its image in $R_{(x)}$ is a nonunit, the polynomial
component $b_R$ lies in $(x)$, whereas
$s_R(0,y)\ne0$.  Choose $a\in k$ with
$s_R(0,a)\ne0$.  The $a$th summand of $N$ then survives
localization at $s$, and any nonzero element $e_a$ of that
summand is annihilated by $b/s^n$.
The same argument works on every smaller principal
neighborhood of $\mathfrak p$.  Such a germ therefore cannot
belong to the stalk of the locally regular sections.

It follows that
$\mathcal M_{X,\mathfrak p}=R_{(x)}$, strictly smaller than
its total quotient field $k(x,y)$; for example $1/x$ belongs
only to the latter.  This also explains why a nonzerodivisor
of one meromorphic stalk need not be a unit of that stalk:
the correct assertion in the next proposition concerns
\emph{locally} regular sections, not all nonzerodivisors
of an individual stalk.
\end{proof}

\begin{proposition}[Proved with corrected definitions: IV 20.1.4, 20.1.8, 20.1.11 and 21.1.5--21.1.6]\label{SA23-D-14}
Define
$\mathcal M_X=\mathcal O_X[\mathscr S(\mathcal O_X)^{-1}]$
by sheafifying sectionwise localization.  Then
$\mathcal O_X\hookrightarrow\mathcal M_X$ is injective and
\[
 \mathscr S(\mathcal M_X)=\mathcal M_X^\times,\qquad
 \mathcal O_X\cap\mathcal M_X^\times
 =\mathscr S(\mathcal O_X).
\]
Thus the total-quotient-ring wording in IV 20.1.4 must be
replaced by localization at the indicated locally regular
denominators; the conclusion of IV 20.1.8 remains true.
For a morphism of ringed spaces $f:X'\to X$, the allowed
denominators in IV 20.1.11 are the sections of
$\mathscr S(\mathcal O_X)$ whose pullbacks are sections of
$\mathscr S(\mathcal O_{X'})$.  They form a subsheaf
$\mathscr S_f$, and localization gives the canonical map of
sheaves of rings on $X'$
\[
 f^{-1}\bigl(\mathcal O_X[\mathscr S_f^{-1}]\bigr)
 \longrightarrow\mathcal M_{X'}.
\]
It extends the structure map
$f^{-1}\mathcal O_X\to\mathcal O_{X'}\hookrightarrow\mathcal M_{X'}$.

For Cartier divisors, the positive cone is the image of
$\mathscr S(\mathcal O_X)$ in
$\mathcal M_X^\times/\mathcal O_X^\times$.  If
$D|_U=\operatorname{div}(g)$ and $f$ is a meromorphic unit,
then $\operatorname{div}(f)|_U\geq D|_U$ exactly when
$f|_U=tg$ for
$t\in\mathscr S(\mathcal O_X)(U)$.
These give the precise replacements in IV 21.1.5 and 21.1.6.
\end{proposition}
\begin{proof}
Every allowed denominator is a nonzerodivisor at every
stalk, hence localization is injective at every stalk
and therefore injective as a sheaf map.
Let $u$ be a locally regular section of $\mathcal M_X$.
Locally write $u=a/t$ with
$a\in\mathcal O_X$ and $t\in\mathscr S(\mathcal O_X)$.
At each point, an equation $a_xb_x=0$ in $\mathcal O_{X,x}$
implies $u_xb_x=0$ in $\mathcal M_{X,x}$; regularity of
$u_x$ and the injection just proved give $b_x=0$.
Thus $a$ is locally regular, and $t/a$ is an inverse
of $u$.  Local inverses agree and glue.
Conversely, a unit is a nonzerodivisor at every stalk.
The same argument with $t=1$ proves the second equality.

Pullback of sections commutes with restriction and gluing,
and local regularity is a local condition.  Thus
$\mathscr S_f$ is a multiplicative subsheaf.
The universal property of localization applies because its
images become units in $\mathcal M_{X'}$.
It supplies the indicated morphism and its uniqueness.
Using merely global nonzerodivisors would not give these
restriction and gluing properties, as
Proposition~\ref{SA23-D-12} shows.

The equality
$\mathscr S(\mathcal O_X)\cap
\mathscr S(\mathcal O_X)^{-1}=\mathcal O_X^\times$
inside $\mathcal M_X^\times$ shows that this positive
cone intersects its negative only in zero, so it defines
the usual order on Cartier divisors.
Finally $\operatorname{div}(f)-D$ is represented on $U$
by the meromorphic unit $f/g$.  Its positivity says
locally that this unit is a section of
$\mathscr S(\mathcal O_X)$; since that is a sheaf, this
is exactly the existence of the asserted section
$t=f/g$ on all of $U$.
\end{proof}

\begin{proposition}[Refuted local formula; proved canonical section: IV 21.2.9]\label{SA23-D-15}
Let $\mathcal J\subset\mathcal M_X$ be an invertible
fractional ideal.  The canonical isomorphism
\[
 \mu:\mathcal M_X(\mathcal J)
 =\mathcal J\otimes_{\mathcal O_X}\mathcal M_X
 \longrightarrow\mathcal M_X
\]
is multiplication.  On an open where
$\mathcal J=\mathcal O_X f$, its coefficient form is
$a\mapsto fa$, not $a\mapsto f^{-1}a$.
For a Cartier divisor $D$, the meromorphic section
$s_D=\mu^{-1}(1)$ of $\mathcal O_X(D)$ satisfies
$\operatorname{div}(s_D)=D$.
Thus L\^e's correction repairs the local identification
and the explanation, while the final displayed conclusion
(21.2.9.3) is correct.
\end{proposition}
\begin{proof}
By definition the map induced by
$\mathcal J\hookrightarrow\mathcal M_X$ sends
$j\otimes a$ to $ja$.
If $\mathcal J=\mathcal O_X f$, write a meromorphic
section in the basis $f$ as $f\otimes a$.
Its image is $fa$.  This is an isomorphism because
$f$ is a meromorphic unit.
The printed reciprocal formula fails, for example, on
$X=\operatorname{Spec}\mathbf Q[z]$ and
$\mathcal J=z\mathcal O_X$: the coefficient $1$
maps by multiplication to $z$, not $1/z$.

Now locally write $D=\operatorname{div}(g)$, so that
$\mathcal O_X(D)=\mathcal O_Xg^{-1}$.
In the basis $g^{-1}$ the section $\mu^{-1}(1)$
has coefficient $g$, since $g^{-1}g=1$.
The divisor of a meromorphic section of an invertible
module is the divisor of its coefficient in a local
regular basis.  Therefore its restriction is
$\operatorname{div}(g)=D|_U$.  These identities glue.
This distinguishes the ambient meromorphic function $1$
from the coefficient $g$ in the chosen line-bundle basis,
which is the missing distinction in the printed explanation.
\end{proof}

\begin{proposition}[Proved: the notation correction in IV 21.6.2]\label{SA23-D-16}
For a locally Noetherian scheme $X$, let
$\mathfrak Z(X)$ be the group of locally finite integer
combinations of irreducible closed subsets, and let
$\mathfrak Z^p(X)$ consist of those supported in pure
codimension $p$ (including the zero cycle).
If $X$ is Noetherian, then
\[
 \mathfrak Z(X)=\bigoplus_{p\geq0}\mathfrak Z^p(X).
\]
Accordingly the final sentence of IV 21.6.2 must refer
to the previously defined $\mathfrak Z^p(X)$, rather than
introducing $Z^p(X)$ with a different typeface.
This is a notation correction, not a false cycle theorem.
\end{proposition}
\begin{proof}
On a quasi-compact space a locally finite family of
nonempty subsets is finite: take finitely many
neighborhoods, each meeting only finitely many members,
that cover the space.  Hence every cycle on Noetherian
$X$ has finite support.  Group its finitely many prime
cycles by the codimension of their generic points.
Different codimensions have disjoint sets of prime
cycles, so this decomposition is unique.  This proves
the displayed direct sum and identifies each graded
piece with exactly the group defined immediately
before it in IV 21.6.2.
\end{proof}

\subsubsection{The literal chapter-zero numbering}

For completeness, a literal reading of the task's
$0_{\rm IV}$ 19.5.5 is covered separately here.
It is the true infinitesimal-testing remark on printed p.~94,
not the graded regular-sequence assertion treated in
Proposition~\ref{SA23-D-5}.
The literal $0_{\rm IV}$ 19.5.6 is headed
\emph{D\'emonstration du th\'eor\`eme (19.5.3)} and has
subparagraphs 19.5.6.1--19.5.6.8; it contains no standalone
assertion numbered 19.5.6(ii).
Its paragraph 19.5.6.6 proves 19.5.3(ii), the
formal-bimorphism statement for symmetric powers of the conormal.
That is a different statement from the false
\emph{IV} 19.5.6(ii) treated in Example~\ref{SA23-D-6}.
\cite[$0_{\rm IV}$ 19.5.3--19.5.6, pp.~91--99]{algEGAIV1}

\begin{proposition}[Proved: the actual $0_{\rm IV}$ 19.5.5]\label{SA23-D-17}
Let $A$ be a topological ring and $B$ a topological $A$-algebra
with its $J$-preadic topology.  Suppose $J$ is maximal,
$C=B/J$ is formally smooth over $A$ for its discrete topology,
and $J/J^2$ is finite dimensional over $C$.
Then $B$ is formally smooth over $A$ if and only if the following
test holds: for every polynomial ring $E=C[T_1,\ldots,T_r]$,
its augmentation ideal $\mathfrak n=(T_1,\ldots,T_r)$, every
$q\geq1$, every
$\mathfrak n^q\subset\mathfrak b\subset\mathfrak n$, and every
augmented $A$-algebra homomorphism $B\to E/\mathfrak b$,
there is an $A$-algebra lift $B\to E/\mathfrak n^q$.
No separatedness of $B$ is required.
\end{proposition}
\begin{proof}
Formal smoothness implies the test, since the displayed
surjection has nilpotent kernel and the augmented map is
continuous: it sends $J^q$ to zero.

Conversely assume the test.  Write $e=\dim_C J/J^2$ and fix
$q\geq2$.  Formal smoothness of the discrete $A$-algebra $C$
splits $B/J^q\to C$ as an $A$-algebra map.  Choose lifts in
$J/J^q$ of a basis of $J/J^2$.  They define a filtered
surjection
\[
 f:C[T_1,\ldots,T_e]/\mathfrak n^q
       \longrightarrow B/J^q
\]
which is the identity on degree zero and an isomorphism
on degree one.  Surjectivity follows by induction on the
finite filtration, since the associated graded of an ideal
is generated in degree one.  Put
$\mathfrak b/\mathfrak n^q=\ker f$.
The test lifts the quotient map $B\to B/J^q=E/\mathfrak b$
to an augmented map $w:B\to E/\mathfrak n^q$.
It kills $J^q$, so it induces
$h:B/J^q\to E/\mathfrak n^q$ with $fh=1$.
On associated graded, $hf$ is the identity in degrees zero
and one, hence in every degree.  Therefore
$\operatorname{gr}f$ is injective as well as surjective.
As $q$ was arbitrary, the canonical map
\[
 \operatorname{Sym}_C(J/J^2)\longrightarrow
                    \operatorname{gr}_J B
\]
is an isomorphism.  In particular every map $f$ just
constructed is an isomorphism.

Here is a direct deduction of formal smoothness, so no
unstated converse criterion is needed.  Let $D$ be a discrete
$A$-algebra, $N\subset D$ a square-zero ideal, and
$u:B\to D/N$ a continuous $A$-map.  Choose $m\geq1$ with
$u(J^m)=0$ and construct the preceding isomorphism for
$q=2m$.  The composite $C\to B/J^{2m}\to D/N$ lifts to
$C\to D$ by formal smoothness of $C$.  Lift the images of
the finitely many $T_i$ arbitrarily to $D$, thereby obtaining
an $A$-map $E\to D$.  Every product of $m$ such images
belongs to $N$, because $u(J^m)=0$.  Every product of
$2m$ belongs to $N^2=0$.  Thus this map kills
$\mathfrak n^{2m}$ and yields the required continuous lift
\[
 B\longrightarrow B/J^{2m}
       \xrightarrow[\sim]{\,f^{-1}\,}
       E/\mathfrak n^{2m}\longrightarrow D.
\]
This proves the remark in precisely its original
finite-conormal scope.
\end{proof}
